\section{A numerical study of continued smoothness}
\label{sec:numerical-campaign}

The first computational question should be precise: for specified smooth
initial packets, how far can solutions of \eqref{eq:baseline} be followed with
resolved derivatives and reproducible error estimates?  A successful finite
interval supplies evidence of continued smooth behavior on that interval.
Extending the interval does not by itself establish global existence.

\subsection{What can already be reused}

The local \texttt{dysthe-core} reference evolves the approved spatial equation
on a periodic profile interval of length $L$, with Fourier differentiation,
retained modes $|m|<N/4$, and integrating-factor fourth-order Runge--Kutta
steps in $\xi$.  Its independent controls include the linear symbol, a nonzero
wavenumber plane wave, the exact NLS traveling profile, and a direct Fourier
convolution of the nonlinear terms including mean flow.  The local
\texttt{dysthe-experiments} code already supplies matched controls,
phase/translation alignment, full periodic residual checks, and separate
grid, step, and domain comparisons.

Stored results for $\epsilon=.025$, $L=160$, $N=4096$,
$h=\Delta\xi=.00125$, and $0\leq\xi\leq4$ give the following maximum raw
complex relative $L^2$ discrepancies.  Each input was compared separately
with $2N$, with $h/2$, and with $2L,2N$ at fixed profile spacing.
\begin{center}
\small
\begin{tabular}{lrrr}
\toprule
Initial field & Grid & Step & Domain\\
\midrule
Solitary profile & $4.17\times10^{-15}$ & $3.07\times10^{-9}$ & $1.54\times10^{-6}$\\
Localized $2\%$ perturbation & $3.09\times10^{-15}$ & $4.08\times10^{-9}$ & $2.26\times10^{-6}$\\
Complex Gaussian & $9.28\times10^{-16}$ & $4.05\times10^{-11}$ & $2.18\times10^{-5}$\\
\bottomrule
\end{tabular}
\end{center}
These are existing two-level agreement checks, not rigorous error bounds or
observed convergence orders.  The stored qualification files were inspected
for this review; these trajectories were not rerun.  Fresh unit tests passed
21 of 21 core checks and 43 of 47 experiment checks; four HTTP tests could not
open a socket in the execution sandbox.  Passing tests neither qualifies a
new packet nor proves a PDE theorem.  The reference additions are local
uncommitted work; a reproducible campaign must record their source hashes.

\subsection{A bounded campaign with explicit stopping rules}

Use $\epsilon\in\{0,.025,.05\}$, treating zero as the mathematical NLS
control.  Start with the existing smooth Gaussian and qualified traveling
profiles, then add signed amplitude and phase perturbations outside the
profile-construction symmetry.  Keep all parameters, initial Fourier
coefficients, and initial-condition identifiers fixed within a refinement
comparison.  The first task is a few documented cases, not a large training
dataset.

Start at $(N,L,h)=(4096,160,.00125)$.  Compare $N,2N,4N$ at fixed $L$ with a
separately qualified step; compare $h,h/2,h/4$ at fixed fine grid; enlarge to
$2L$ and, when necessary, $4L$ at fixed profile spacing.  Reconstruct localized
initial data on each domain rather than repeating its periodic copies.
Recomputed solitary profiles must retain and report their initial-domain
discrepancy.  The nonlocal mean flow can change with $L$ even when edge
amplitudes are tiny.

Use successive horizons
\[
  \Xi\in\{1,2,4,8,16,32\},
\]
advancing only after the preceding horizon passes its checks.  Compare all
saved stations and inspect denser output near rapid growth.  Retain raw
complex errors as well as errors after fitting only translation and constant
phase; never fit away amplitude changes.  Perturbations should also be
compared with their matched propagated control.  A failed gate calls for
refinement or a shorter qualified interval; it is not evidence of blow-up
or nonexistence.

For $k_m=2\pi m/L$ and Fourier coefficients normalized by $1/L$, record
\begin{equation}\label{eq:monitor-sobolev}
 \norm{u}_{H^s(\TT)}^2
 =L\sum_m(1+k_m^2)^s|\widehat u_m|^2,
 \qquad s=1,2,3.
\end{equation}
Also record $\norm{u}_{\infty}$, $\norm{u_\tau}_{\infty}$, and spectral
slopes above the resolved signal floor.  Monitor the fraction of each
$H^s$ sum in the upper fifth of the retained band.  A small unweighted
spectral tail can conceal a large derivative error.  Check boundary and
exterior contributions in derivative norms, not action alone.  Apparent
norm growth is interpretable only while these weighted tails, slopes, and
refinement comparisons remain resolved.

Provisional gates are maximum raw relative $L^2$ discrepancy $10^{-4}$,
$H^1$ discrepancy $10^{-3}$, and $H^2,H^3$ discrepancies $10^{-2}$, with
upper-band weighted fractions below $10^{-6}$.  Use the finer field for
normalization and absolute scales when the reference norm is near zero.
For a weak perturbation, additionally require the signed response
discrepancy below $2\%$ of the measured response above its numerical floor.
These are proposed acceptance tolerances, not established error bounds;
tighten them when the scientific effect is smaller.  Step halving must show
the expected order before roundoff, rather than just two close answers.
The existing $\epsilon\norm{u}_\infty$ and bandwidth screens remain
separate physical-model checks; their thresholds are not universal validity
theorems.

Evaluate the unprojected RHS on an enlarged Fourier grid, retaining discarded
nonlinear modes.  For a continuous reconstruction $v$ of a numerical
trajectory, the useful defect is $r=v_\xi-F(v)$, including error from the
reconstruction in $\xi$.  The existing snapshot RHS diagnostic alone does
not provide that time-dependent defect or its integral.  Mass drift is an
additional diagnostic; it cannot replace derivative resolution or residual
control.

\section{A Fourier-in-profile physics-informed network}
\label{sec:spectral-pinn}

\subsection{Represent the periodic field and retain the full residual}

The recommended first PINN is a neural approximation of the Fourier
coefficients as functions of propagation distance:
\begin{equation}\label{eq:spectral-pinn}
 v_\vartheta(\xi,\tau)
 =\sum_{m=-K}^{K}a_m(\xi;\vartheta)e^{ik_m\tau},
 \qquad a_m\in\CC.
\end{equation}
A smooth network takes $\xi$ and returns the real and imaginary parts of
all coefficients.  Periodicity is exact.  Because the envelope is complex,
\emph{do not impose} $a_{-m}=\overline{a_m}$ or conjugate-reflection symmetry.
Choose $K$ from convergence of the reference fields in $H^3$, and repeat
with a larger $K$; action capture alone is insufficient.

Automatic differentiation supplies $a_m'$, while the profile derivatives
use $ik_m$ and $-k_m^2$.  At each sampled $\xi$, form
$\rho=|v_\vartheta|^2$, compute $\D\rho$ with multiplier $|k_m|$ and zero
mean mode, and evaluate every term of \eqref{eq:baseline}.  Define
\begin{equation}\label{eq:pinn-residual}
 r_\vartheta=\partial_\xi v_\vartheta-F(v_\vartheta).
\end{equation}
The cubic terms, including $v_\vartheta\D(|v_\vartheta|^2)$, can contain modes
through $3K$.  The residual must retain these modes even though the trial
field stops at $K$.  A uniform quadrature grid with $Q>6K$ points resolves
this full residual and its squared $L^2$ norm for the trigonometric trial
field.  Penalizing only $|m|\leq K$ could conceal a large discarded RHS and
merely solve a truncated coefficient system.  Independently validate with
larger $Q$, larger $K$, and new propagation stations.

\subsection{Hard initial values and short propagation windows}

On $[\xi_j,\xi_{j+1}]$, write
\[
 a_m(\xi;\vartheta)
 =a_m^{(j)}+(\xi-\xi_j)b_m(\xi;\vartheta).
\]
This enforces the supplied initial coefficients exactly.  On the first
window, report the projection error of the prescribed initial field onto
$|m|\leq K$; a hard condition does not remove this error.  On later windows,
set $a_m^{(j)}$ to the previous learned terminal value.  Errors therefore
accumulate honestly.  Replacing every window's initial value with a
reference snapshot defines a different, data-assisted experiment and must
be reported separately.

Begin with short windows inside $\Xi=1$, double precision, a smooth modest
network, and at least three recorded initialization seeds.  A practical
training sequence is Adam followed by a bounded L-BFGS refinement.  Declare
the architecture, coefficient scaling, optimization budget, window length,
and stopping rules before comparison.  A basic loss is
\[
 \mathcal L_{\rm PDE}
 =\frac{1}{(\xi_{j+1}-\xi_j)L}
   \int_{\xi_j}^{\xi_{j+1}}\int_{\TT}|r_\vartheta|^2\dd\tau\dd\xi.
\]
An explicitly stated Fourier-weighted residual may improve derivative
control, but changes the training objective and must be tested.  Conservation
is a diagnostic rather than a substitute for the PDE loss.  Do not supply
interior reference snapshots to a purported physics-only run; if such data
are used, report that variant alongside it.

\subsection{What would count as a meaningful result}

First require the NLS and plane-wave controls, then the nonconstant mean-flow
case, then the packet campaign.  At held-out stations compare the network
with the refined reference in raw complex $L^2$, $H^1,H^2,H^3$, peak and
profile slope, phase where amplitude is meaningful, and full residual.
Compare matched perturbation responses before interpreting their growth.
Training failure, a small collocation loss, or disagreement between two
networks cannot decide existence.  A residual becomes an error certificate
only when combined with justified stability estimates or validated inverse
bounds of the kind discussed elsewhere in this report.

The active learning repository has no implemented water-wave PINN.  This
review adds no new trajectory or trained network.  Its proposed first
deliverable is one reproducible short-window comparison, followed by horizon
extension only after qualification.  Record total training and evaluation
time, peak memory, all restarts and seeds, and any reference-generation cost;
compare with the numerical solver at matched accuracy.  A PINN has no
presumed speed advantage for this one-coordinate problem.  Its potential
value must come from a demonstrated use---for example, a differentiable
family of accurate approximations---rather than the name of the method.
